% Limitations and Scope — Securing LLMs Against Prompt Injection
% Academic manuscript section (IEEE / arXiv style)

\section{Limitations and Scope}
\label{sec:limitations}

This section clarifies what the present study does and does not claim.
We first state the intended scope of the threat model, data, and evaluation protocol, then enumerate limitations that qualify the empirical conclusions.

\subsection{Scope}
\label{sec:scope}

The study addresses \emph{direct} textual prompt injection against an instruction-tuned chat model under a model-agnostic pre-inference gate.
The defender may inspect and optionally block the user prompt before generation, but may not retrain the backend model or redesign the application interface into structured queries~\cite{struq2024,secalign2024}.
Within that setting we compare lexical baselines with two context-aware sanitizers: hybrid Method~A, in which semantic hard triggers form the primary layer and weighted multi-signal scoring provides secondary coverage, and Method~B, which relies on learned soft fusion without hard vetoes.

The evaluation corpus is intentionally compact and curated for controlled comparison.
Primary metrics are reported on a frozen held-out test partition using a two-layer protocol that separates Bypass Rate from True ASR, complemented by paraphrase, edge-case benign, adaptive rewriting, ablation, and latency analyses.
The contribution is therefore a reproducible proof-of-concept account of where defensive capability originates in a lightweight sanitization pipeline, not a production security certification.

Out of scope for the present experiments are agentic tool-use settings~\cite{debenedetti2024agentdojo}, long-context indirect injection through retrieved documents or web content as the primary attack channel~\cite{greshake2023indirect,hines2024spotlighting}, multimodal injection~\cite{yeo2025multimodal}, and unbounded optimization-based jailbreak search~\cite{yi2024jailbreak}.
Training-time alignment and structured-query redesign are treated as complementary defenses rather than competing systems under evaluation.

\subsection{Limitations}
\label{sec:limitations-detail}

Several limitations qualify the reported results.

\paragraph{Data scale and statistical power.}
The corpus contains 350 prompts with a 53-prompt test partition.
On a set of this size, leave-one-component ablations that show zero change should not be over-interpreted as proof that a signal is universally useless; they indicate only that the component did not alter decisions on this particular hold-out when the remaining cues were present.
Larger and more diverse benchmarks would be needed for confident estimates of ablation deltas and confidence intervals.

\paragraph{Threat-model coverage.}
We evaluate direct single-turn textual injection.
Indirect poisoning, multi-turn manipulation, retrieval-mediated attacks, and tool-diverting agent attacks may expose failure modes that the present suite does not stress~\cite{greshake2023indirect,debenedetti2024agentdojo,wu2025epistemic}.
Adaptive rewriting uses a compact sequence-to-sequence model and therefore probes controlled style transfer rather than the strongest available adversarial optimizers~\cite{yi2024jailbreak,zizzo2025adversarial}.

\paragraph{Judge and target dependence.}
True ASR depends on an LLM-as-a-judge from the same model family as the target.
Under greedy decoding the protocol is reproducible for a fixed model revision, but the judge is not a human gold standard and may disagree with other models or annotators~\cite{zizzo2025adversarial}.
Results may also shift if a stronger or differently aligned backend replaces Qwen2.5-3B-Instruct.

\paragraph{Architectural interpretation.}
Method~A is a hybrid system: ablations show that semantic hard triggers provide most of the observed protection, while soft weighted scoring is complementary.
Claims that would present Method~A as a balanced multi-signal ensemble would overstate the role of soft aggregation on this benchmark.
Method~B improves over lexical baselines with milder false positives, but does not match Method~A's static-test True ASR; soft fusion alone is therefore not a drop-in substitute for the hybrid gate under the present evaluation.

\paragraph{Usability cost.}
Method~A's stronger security coincides with higher false-positive rates on ordinary benign prompts and especially on technical edge-case prompts.
This trade-off is intrinsic to the current operating point and may limit deployment in product settings where over-blocking is costly~\cite{bair2025prompt,kholkar2025capture}.
Threshold retuning or benign whitelisting could reduce that cost, but such changes were not explored here in order to keep the architecture and evaluation fixed.

\paragraph{Efficiency and deployment.}
Sanitizer auxiliaries run on CPU by design; Method~A adds tens of milliseconds of pre-inference latency relative to lexical filters.
This overhead is modest compared with large-model generation, but it is not free, and wall-clock cost will vary with hardware and batching.
The study does not evaluate multi-tenant serving, streaming interfaces, or privacy constraints on prompt logging.

Taken together, these limitations imply that the results should be read as evidence about hybrid versus soft context-aware sanitization under a constrained direct-injection protocol.
They support retaining and transparently reporting the present architecture, while leaving broader generalization, adaptive weighting, and reduced hard-trigger reliance to future work.
